% Separate post-CP29 candidate; no unrestricted first-wild least-index claim.
\subsection{Two translation breaks recover a tame residual Newton parameter}

Let \(p\ge5\), let \(K\) be complete discretely valued of mixed
characteristic \((0,p)\) with algebraically closed residue field,
and suppose \(\mu_p\subset K\). Its absolute index \(e\) is
arbitrary; in particular it may be divisible by \(p\).
The cyclotomic uniformizer valuation gives \(p-1\mid e\).
Put \(q=p^2\), \(N=q-1\), \(s=(p+1)e\), and pass to
the common tame good field \(L/K\) of degree \(N\).
The groups considered are whole potentially supersingular
height-two groups.

\begin{proposition}[One prime parameter suffices in the tame residual comparison]
\label{prop:B-RD84-tame}
\status{Candidate; full proof supplied, independent audits pending.}
Suppose two such groups have two-slope first-level polygons,
\(L\)-rational first-level modules, and isomorphic actual
\(p^2\)-torsion. Write their parameters and valuations as
\[
 \begin{gathered}
 B_i=pb_i,\quad0<b_i<e,\qquad
 \alpha_i=s-pb_i,\qquad\beta_i=b_i,\\
 v_L(T_i(Q))=\alpha_i\quad(0\ne Q\in H_i),\qquad
 v_L(T_i(Q))=b_i\quad(Q\notin H_i).
 \end{gathered}                                      \tag{B-RD84.1}
\]
If at least one of \(b_1,b_2\) is prime to \(p\), then
the actual comparison identifies the high lines,
\(b_1=b_2\), and the finite special characters agree.
No bound on the full second- or third-level division-field
degree is assumed.
\end{proposition}
\begin{proof}
Index the models so that \(p\nmid b_1\). Choose corresponding
residual vectors outside both high lines and corresponding
primitive second-level lifts \(P_i\). The actual inverse
Newton polygon gives
\[
                     v_L(T_i(P_i))=b_i/q.             \tag{B-RD84.2}
\]
Here the same universal coefficient and line-integrality
proof as in B69 and B82 applies at arbitrary \(e\);
no coefficient argument requires \(p\nmid e\).

The actual torsion isomorphism identifies the point
stabilizers, hence the point fields \(E=L(P_1)=L(P_2)\).
Because both first-level modules and both \(pP_i\) are
\(L\)-rational, each field is Galois with its group a
subgroup of the first-level translation group. Indeed
every conjugate is \(P_i+Q\) for an \(L\)-rational \(Q\),
and so is defined over the same field. Its degree is at
most \(q\). The reduced denominator of \(b_1/q\) is
\(q\), forcing \([E:L]=q\). Thus every first-level
translation occurs in the same actual point field, even
if \(b_2\) is divisible by \(p\).

Normalize \(v_E\) to have value group \(\mathbf Z\).
The coordinates of \(P_i\) have valuations \(b_i\).
For every nonidentity \(g\in\operatorname{Gal}(E/L)\),
Lemma~\ref{lem:B-RD60-break} gives the exact equality
for model 1 and the unconditional inequality for model 2:
\[
 \begin{split}
 v_E(gT_1(P_1)-T_1(P_1))&=b_1+r(g),\\
 v_E(gT_2(P_2)-T_2(P_2))&\ge b_2+r(g).
 \end{split}                                         \tag{B-RD84.3}
\]
Both concern the same actual lower break \(r(g)\).

If the high lines differed, choose a translation
\(0\ne Q_1\in H_1\) corresponding to \(Q_2\notin H_2\).
It occurs because the full first-level translation group
occurs. Its high displacement in model 1 gives
\(r(g)=q\alpha_1-b_1\), whereas its low displacement
in model 2 would require
\[
                  qb_2\ge b_2+q(s-pb_1)-b_1.         \tag{B-RD84.4}
\]
But \(\alpha_1=s-pb_1>e>b_1,b_2\), contradicting this
inequality. Explicitly the difference of the left side
and right side is increasing in both integers and at
\(b_1=b_2=e-1\) equals \(-q(p+1)<0\).
The high lines therefore agree.

Choose a nonzero low translation in the common module.
Its break from model 1 is \((q-1)b_1\); (B-RD84.3)
in model 2 gives
\[
                     (q-1)(b_2-b_1)\ge0.             \tag{B-RD84.5}
\]
Choose instead a nonzero high translation. Its break is
\(q(s-pb_1)-b_1\), and the other displacement gives
\[
                     (p^3+1)(b_1-b_2)\ge0.           \tag{B-RD84.6}
\]
These force \(b_1=b_2\), with no prime-to-\(p\) assumption
on the second coordinate. The identified high-line
characters are \(\zeta^{c_i-B_i}\), in an exact tame
Honda coordinate. Since \(B_1=B_2\), their equality
gives \(c_1=c_2\pmod N\), and hence equal finite
special characters.
\end{proof}

\begin{corollary}[The remaining tame reducible first-wild parameters]
\label{cor:B-RD84-first-wild}
\status{Candidate; full proof supplied, independent audits pending.}
At \(e=p(p-1)\), a different-character actual
\(p^3\)-congruent pair with tame reducible residual inertia
must have two-slope parameters
\[
 B_i=p^2k_i,\qquad1\le k_i\le p-2,                  \tag{B-RD84.7}
\]
and its high lines are exchanged. This restriction is
independent of the size of its full wild torsion image.

At \(p=5,e=20\), after a common finite height-one twist
normalizing the nonsquare finite determinant, the unordered
parameter pair is exactly \(\{25,75\}\). Its finite
types are A/B up to common quadratic twist, with A, B,
C as in B77. This is only a necessary pattern and asserts
no existence of such a pair.
\end{corollary}
\begin{proof}
The square determinant case is already excluded at the
third level by B77. A one-slope residual representation,
when tame and reducible, acts by a scalar on its
two-dimensional leading-coordinate plane and has square
determinant. Thus both polygons have two slopes. Since
their residual action is tame, both first-level modules
are rational over \(L\), and \(B_i=pb_i\).

Their actual \(p^2\)-torsion is already isomorphic.
The proposition therefore forces both \(b_i\) to be
divisible by \(p\) if the finite characters differ.
The range \(0<b_i<e=p(p-1)\) gives (B-RD84.7).
If their high lines agreed, equality of finite determinants
would give \(c_1\equiv c_2\pmod{p-1}\), and the common
high-line character would give
\(p^2k_1\equiv p^2k_2\pmod{p-1}\). The range
\(1\le k_i\le p-2\) forces \(k_1=k_2\); the same
high-line character then gives \(c_1=c_2\pmod N\).
This is excluded, so the high lines must be exchanged.

For the explicit pair at \(p=5\), make the common
nonsquare determinant primitive by the common finite
twist already justified in B50 and B77. This leaves the
good models over \(L\), hence their \(B_i\), unchanged.
Now \(c_i\equiv1\pmod4\), while
\(B_i\equiv c_i\pmod6\). Thus \(k_i\) is odd in
\(\{1,2,3\}\), giving \(B_i\in\{25,75\}\).
Equality of the residual characters with exchanged high
lines forces the two distinct values. More explicitly,
the possible normalized ordered pairs are
\[
 (B_1,c_1;B_2,c_2)=(25,1;75,9)
             \quad\text{or}\quad(25,13;75,21),         \tag{B-RD84.8}
\]
up to interchanging the two models. Their residual
high/low exponents are respectively \((0,6;6,0)\)
and \((12,18;18,12)\) modulo 24. The finite special
orbits are A/B and C/B; a common quadratic twist
interchanges A and C and preserves B. This proves the
stated necessary pattern directly.
\end{proof}

\begin{remark}[What this parameter recovery does not decide]
\label{rem:B-RD84-scope}
\status{Candidate scope; the surviving exchanged-line pattern remains open.}
The proposition uses the small common point field of a
single primitive second-level point. It does not assume
that the full division field has that degree. Both the
high and the low translations in that same field are
essential; using only the low inequality would leave
spurious one-divisible parameter candidates.
The corollary removes the earlier exploratory tame pairs
\(\{35,75\}\) and \(\{55,75\}\) at every image size.
The pair \(\{25,75\}\), all larger-image questions in
other residual sectors, and the full index interval
\(\{2,3\}\) remain open. No formal-group pair is
constructed by the finite parameter enumeration.
\smallskip
\noindent\emph{Exact supplement and provenance.}
The retained program
\path{code/reduction_descent/verify_tame_residual_one_prime_recovery.py}
checks the three symbolic displacement identities,
80,536 point-field parameter rows over 96 prime/index
samples, including 6,600 mixed prime/divisible exclusions,
and the exact two normalized \(p=5\) character rows and
their common quadratic-twist orbit. All parameter pairs
are checked for \(p=5,7,11,13\) at the four displayed
index factors; larger primes through 101 use explicitly
retained boundary samples. The author run passed in
\(0.12348706599732395\) seconds with no failed execution.
The uniform common-point-field and two-break proof is
independent of those finite ranges. All executed code,
output and prior source versions remain retained.
The exact internal B69 quotient and inverse-polygon proof,
B60 fixed-coefficient break lemma and B82/B83 common-field
arguments were reread; no new original-source read is
claimed.
\end{remark}
